\section{Diskussion}
\label{sec:Diskussion}
Alle Experimentteile werden aus mehreren Metallteilen zusammengesetzt, deren Kanten sowohl innen als auch außen entgratet sind.
Es ist daher möglich, dass Abweichungen in den Messwerten durch Randeffekte an den Enden der Zylinder oder dem Rand der Kugelhälften entstehen.
Da alle Bauteile eine ebene Oberfläche haben, sollten die Abweichungen jedoch gering sein.
\subsection{Messung der Schallgeschwindigkeit}
Die beiden unabhängig voneinander ermittelten Schallgeschwindigkeiten $c_{50}$ und $c_{75}$ weichen zueinander um $\increment c_{\exp} = \qty{0.85(0.37)}{\percent}$ ab.
Diese gute Übereinstimmung zeigt, dass die Art des verwendeten Aufbaus nur wenig Einfluss auf die Ergebnisse hat.
Insbesondere ist die höhere Anzahl an Grenzflächen in dem Versuchsteil mit $50\,\unit{\milli\meter}$-Zylindern keine signifikante Fehlerquelle.
Aufgrund der geringen Abweichung der beiden Werte zueinander ist der berechnete Mittelwert eine gute Beschreibung der Schallgeschwindigkeit.

\noindent
Der Literaturwert der Schallgeschwindigkeit in trockener Luft bei $T = \qty{20}{\celsius}$ ist $c_{\mathrm{theo}} = \qty{344}{\meter\per\second} $~\cite[22]{friesecke_2014}.
Die Abweichung zum experimentellen Mittelwert beträgt $\increment c_{\mathrm{theo}} = \qty{3.98(0.18)}{\percent}$.

\noindent
Der Raum, in dem das Experiment durchgeführt wurde, ist nicht klimatisiert und hat keine Lufttrocknungsanlage. 
Es ist daher denkbar, dass die Abweichung zwischen Theorie und Experiment durch Abweichungen in Temperatur und Luftfeuchtigkeit zustandekommt.
%
%
%
\subsection{Orbitale des Wasserstoffatoms}
In \autoref{fig:H180_Spektrum} sind ab $\qty{6}{\kilo\hertz}$ sekundäre Resonanzpeaks neben den Hauptresonanzen erkennbar.
Dies deutet auf Überlagerungen der Resonanzen im Kugelresonator hin, welche die Messwerte stören können.

\noindent
\autoref{fig:resonanzen} zeigt, dass die Messwerte in allen vier Fällen erkennbar zu den Theoriefunktionen passen.
Die $\qty{2.31}{\kilo\hertz}$-Resonanz weist Abweichungen von der Theorie um $\vartheta=\qty{90}{\degree}$ auf.
Die Druckamplitude steigt, anstatt auf Null zu fallen.
Bei der $\qty{3.705}{\kilo\hertz}$-Resonanz ist die Abweichung nach oben ebenfalls erkennbar.
Die Form der Messwerte folgt allerdings der Theoriekurve.
Für die beiden höheren Resonanzfrequenzen passen die Messwerte sehr gut zur Theoriekurve.
Die Messwerte nähern sich scheinbar mit höherer Ordnung der Theorie besser an.
Das deutet darauf hin, dass die dominante Mode eine Mode höherer Ordnung ist und die niedrigeren Moden durch Überlagerungseffekte gestört werden.
Die $\qty{2.31}{\kilo\hertz}$-Resonanz entspricht der Quantenzahl $l=1$.
Die nächsten höherfrequenten Resonanzen entsprechen den Quantenzahlen $l=2$, $l=3$ und $l=5$.
$l=4$ wurde nicht gemessen, da zwischen der letzten und vorletzten Resonanz ein Resonanzpeak ausgelassen wurde.
Dies zeigt, dass die Ordnung der Resonanzen, angefangen bei der Grundresonanz $f_0$, der Quantenzahl $l$ entsprechen.

\noindent
Die Vergrößerung der Entartung durch Einsetzen dickerer Metallringe zwischen die Kugelhälften geht aus \autoref{fig:aufspaltung} deutlich hervor.
\autoref{fig:ring_fit} zeigt eine lineare Proportionalität zwischen Ringdicke und Frequenzabstand der entarteten Resonanzen.
Die Fehlerbalken der Messwerte liegen in der Größenordnung $10^1$, was nur eine Potenz unter der Größenordnung der Messwerte liegt.
Der Messfehler ist also groß.
Eine mögliche Fehlerquelle ist, dass der $\qty{12}{\milli\meter}$-Ring aus dem $\qty{3}{\milli\meter}$- und dem $\qty{9}{\milli\meter}$-Ring zusammengesetzt ist.
Wenn die Ringe nicht optimal aufeinander passen, ergibt sich ein systematischer Fehler in der Ringdicke.
Zur Verbesserung des Ergebnisses kann ein nicht-zusammengesetzter Ring verwendet oder der zusammengesetzte Ring mit einem Messschieber ausgemessen werden.

\noindent
Die hochaufgelöste Vermessung der $\qty{9}{\milli\meter}$-Aufspaltung zeigt, dass die Entartung die Quantenzahl $l$ ändert, aber $m$ konstant bleibt.
\autoref{fig:9mm_polar} zeigt, dass die Messwerte für die beiden entarteten Resonanzen sehr gut zu den Theoriekurven der entsprechenden Kugelflächenfuktionen passen.
Dies bestätigt die Annahme der Entartung in $l$.
